\begin{resumen}

    En el siguiente proyecto se desarrolla una solución para extraer, validar y estructurar automáticamente información contenida en documentos financieros en formatos digitales, tales como PDF e imágenes. El trabajo está orientado a entornos \textit{fintech} donde la revisión manual es susceptible a errores humanos, esto demanda tiempo y limita la capacidad operativa del servicio. El desarrollo del proyecto incluye la preparación del \textit{dataset}, el procesamiento/validación de documentos, y la evaluación comparativa de ocho estrategias de extracción organizadas en tres fases: métodos determinísticos basados en expresiones regulares, enfoques híbridos que combinan OCR con modelos de lenguaje, y el envío directo del PDF al modelo. La evaluación se realizó sobre 176 estados de cuenta bancarios reales de 11 instituciones financieras mexicanas, y se determinó que el envío directo de un PDF al modelo de lenguaje ofrece los mejores resultados, alcanzando una precisión promedio de 93.4\% en los tres campos evaluados (92.8\% en Titular, 90.9\% en CLABE y 96.6\% en Banco), con una mediana de procesamiento de 1.9 segundos y un costo de \$0.011 USD por documento. A partir de estos hallazgos, se construyó un producto mínimo viable (\textit{MVP}) basado en extractores configurables mediante \textit{schemas}, \textit{prompts} y modelos, lo que permite adaptar el sistema a distintos tipos de documentos sin modificar el código fuente. La variabilidad de precisión en las pruebas de extracción de datos bancarios fundamentó la decisión de delegar la especialización al usuario mediante herramientas de asistencia basadas en inteligencia artificial. Además, la solución incorpora un módulo que almacena resultados, correcciones manuales y métricas de desempeño, con lo cual es posible consolidar una línea base confiable y sostener la mejora continua de los extractores. De este modo, el \textit{MVP} no solo constituye una implementación funcional para un caso de uso específico, sino también una plataforma reutilizable para otros escenarios de automatización de documentos en entornos \textit{fintech}.
    
    \noindent \textbf{Palabras clave:}\\
    \noindent Extracción de documentos; Modelos de lenguaje; Documentos bancarios; Salida estructurada; \textit{Fintech}
    \end{resumen}
    